% paper/sections/08c_defect_correction.tex
%
% §8.3  欠陥補正法（Defect Correction Method）による PPE 求解
%   — §8.2 で構築した CCD 離散演算子 L_H を高精度・低コストで解く枠組み
%   — Thomas 三重対角係数の導出：付録 \ref{app:sweep_thomas_coeff}
%   — 収束理論・スペクトル解析：付録 \ref{app:dc_convergence_theory}

\subsection{欠陥補正法による効率的 PPE 求解}
\label{sec:defect_correction_main}

\noindent\textbf{本節の位置づけ：}
本節で示す欠陥補正法（DC）は変密度 PPE 求解の基本フレームワークである．
反復回数 $k$ に応じて精度が向上し，$k \geq 3$ で $\Ord{h^6}$ に漸近する
（§~\ref{sec:dc_accuracy}）．$k$ の選択および FD 直接法 + CCD 勾配の設計根拠は
§~\ref{sec:pressure_accuracy_summary} にまとめる．

本節では，§~\ref{sec:ccd_poisson_matrix} で構築した CCD 離散演算子 $L_H$（$\Ord{h^6}$）を
「\textbf{高精度で評価し，低精度で解く}」という原理で効率的に求解する
\textbf{欠陥補正法（Defect Correction Method）}~\cite{Stetter1978} を示す：
残差（欠陥）の評価には $L_H$ を用いて $\Ord{h^6}$ 精度を確保し，
線形系の求解には方向分離が容易な低次演算子 $L_L$（2次精度 FD）を用いて
$\mathcal{O}(N^2)$ コストを維持する．

§~\ref{sec:ccd_poisson_matrix} で CCD 空間離散化スキームを選択済みであり，
本節はその CCD 系を解く\textbf{線形ソルバ}の選択に特化する（ソルバ比較は §~\ref{sec:dc_comparison}）．

\noindent\textbf{なぜ $L_H$ を直接解かないのか：}
$L_H$ は CCD ブロック Thomas 法に基づく密結合演算子であり，
$x$-$y$ 方向の結合が強いため ADI 的な方向分離が成立しない．
大域疎行列の陽的構成とその直接解法は付録~\ref{app:ccd_lu_direct} に示す．

\subsubsection{アルゴリズム}
\label{sec:dc_algorithm}

離散化された PPE $L_H\,p = b$ に対し，
初期推定 $p^{(0)} = 0$ から以下の3ステップを反復する：
\begin{equation}
  \boxed{%
  \begin{aligned}
    &\textbf{Step 1（欠陥の算出）：}\quad
      d^{(k)} = b - L_H\,p^{(k)}
    \\[4pt]
    &\textbf{Step 2（補正の求解）：}\quad
      L_L\,\delta p^{(k+1)} = d^{(k)}
    \\[4pt]
    &\textbf{Step 3（解の更新）：}\quad
      p^{(k+1)} = p^{(k)} + \omega\,\delta p^{(k+1)}
  \end{aligned}
  }
  \label{eq:dc_three_step}
\end{equation}

各記号の定義：
\begin{itemize}
  \item $L_H$：CCD 法で離散化した可変係数 Laplacian（$\Ord{h^6}$，§~\ref{sec:ccd_poisson_matrix}）．
        行列不要で評価可能（各行・列に CCD ブロック Thomas 法を適用）．
  \item $L_L$：2次精度中心差分（FD）で離散化した近似 Laplacian．
        三重対角構造を持ち，$x$・$y$ 方向に\textbf{方向分離（directional sweep）}が可能．
  \item $d^{(k)}$：欠陥（defect）— 現在の近似解 $p^{(k)}$ を
        高精度演算子 $L_H$ で評価した残差．
  \item $\delta p^{(k+1)}$：補正量（correction）— 低次演算子 $L_L$ の逆作用で得る．
  \item $\omega$：緩和係数（relaxation parameter）—
        収束条件 $\omega < 2/r_{\max} = 0.833$ の詳細は付録~\ref{app:dc_convergence_theory}．
\end{itemize}

\subsubsection{Step~2 の求解}
\label{sec:dc_sweep}

Step~2 の補正方程式 $L_L\,\delta p = d$ は，
$L_L$（変密度 FD Laplacian，Neumann BC，ゲージピン固定）を
\textbf{直接法}で厳密に求解する：
\begin{equation}
  \delta p^{(k+1)} = L_L^{-1}\,d^{(k)}
  \label{eq:dc_step2_lu}
\end{equation}
$L_L$ の Neumann BC はゴーストセル反射 $p_{-1}=p_1$（左壁）で離散化し，
$L_H$（CCD）と $L_L$（FD）が同一の物理的 Neumann 条件を共有する
（境界スタンシルの導出は付録~\ref{app:sweep_thomas_coeff}）．
直接法の採用は PPE 残差を機械精度まで消去し，
反復法における許容残差 $\epsilon_\mathrm{tol}$ の系統誤差蓄積を原理的に排除するためである．
方向分離（ADI 型）スイープとの比較解析は §~\ref{sec:verify_dc_spectral} および付録~\ref{app:dc_convergence_theory} を参照．

\subsubsection{収束先の精度保証}
\label{sec:dc_accuracy}

欠陥補正法の本質的な性質は，
\textbf{左辺に低次演算子 $L_L$ を用いても，高精度演算子 $L_H$ の精度が維持される}点にある．
収束条件 $\|d^{(k)}\| < \varepsilon_\mathrm{tol}$ を満たした時点で，
\begin{equation}
  L_H\,p^{(k)} \approx b
  \quad\text{（精度 $\varepsilon_\mathrm{tol}$）}
\end{equation}
が成立し，最終圧力解は CCD 演算子 $L_H$ の $\Ord{h^6}$ 精度を完全に保持する．
$L_L$ は収束を駆動する「加速装置」の役割のみを担い，
収束先の精度は右辺の欠陥評価（$L_H$ による残差）が決定する．

この精度保証は反復回数 $k$ と密接に関連する：
$k=1$ では $L_L$ の $\Ord{h^2}$ が支配的だが，
$k$ を増やすごとに $L_H$ の高次精度が発現し，
$k \geq 3$ で $\Ord{h^6}$ に漸近する
（定量的検証は §~\ref{sec:dc_iteration_accuracy} に示す）．
反復行列のスペクトル解析（付録~\ref{app:dc_convergence_theory}）により，
$\omega$ の最適値と収束速度が理論的に導かれる：
$\omega \in (0,\, 0.833)$ が収束に必要であり，
$\omega \in \{0.3, 0.5, 0.7\}$ はいずれも等密度条件で良好な収束を示す
（実験検証: §\ref{sec:verify_dc_spectral}）．

\phantomsection\label{sec:dc_comparison}% backward-compat

